\begin{frame}[t]{\ev{\tagO\,\tagB\,\tagP}Kernel isolation and authority are separate controls}
\vspace{0pt}
{\small\renewcommand{\arraystretch}{1.2}
\begin{tabularx}{\textwidth}{@{}L{5.3cm} L{3.2cm} Y@{}}
\toprule
Configuration & Kernel separation & Publishing authority\\
\midrule
microVM containing a bearer token & hypervisor boundary & \leavevmode{\color{Rust}a copied token can remain usable}\\
worker, with controller-mediated publication & depends on the runtime & \leavevmode{\color{Teal}the controller checks each publication}\\
\bottomrule
\end{tabularx}\par}
\vspace{.2cm}
{\small
\textbf{Kernel surface:} containers, gVisor, microVMs and unikernels expose different amounts of kernel.\par\vspace{2pt}
\textbf{Explicit authority:} Capsicum and CHERI give a process only the rights it is handed.\par\vspace{4pt}
\tagP\ \textbf{Sharing policy for children}\par\vspace{2pt}
\begin{ticks}
\item warm dependency cache: immutable, identified by digest, shared within one trust domain
\item writable branch state: private to the child
\end{ticks}\par}
{\footnotesize\color{Muted}Immutability does not remove timing or content leakage, so a shared cache still needs a trust policy.\par}
\sources{Firecracker (NSDI'20); gVisor documentation; Watson et al., Capsicum (USENIX Sec'10) and CHERI (IEEE S\&P'15).}
\note{[0:40] Kernel isolation and authority are separate controls. Containers, gVisor, microVMs and unikernels differ in how much kernel the guest can reach. Capsicum and CHERI address explicit authority. A microVM that holds a bearer token has a hypervisor boundary, but a copied token stays usable. A shared cache must be immutable and identified by digest, and it still needs a trust policy. Forking virtual machines is an old idea.}
\end{frame}
